\documentclass[openleft]{memoir}
\setulmarginsandblock{1.5cm}{1.5cm}{*}
\setlrmarginsandblock{2.5cm}{1.5cm}{*}
\checkandfixthelayout
\usepackage{amsmath}
\usepackage{authblk}
\usepackage{listings}
\usepackage{color,calc,graphicx,soul,fourier}
\definecolor{nicered}{rgb}{.647,.129,.149}
\makeatletter
\newlength\dlf@normtxtw
\setlength\dlf@normtxtw{\textwidth}
\def\myhelvetfont{\def\sfdefault{mdput}}
\newsavebox{\feline@chapter}
\newcommand\feline@chapter@marker[1][4cm]{%
\sbox\feline@chapter{%
\resizebox{!}{#1}{\fboxsep=1pt%
\colorbox{nicered}{\color{white}\bfseries\sffamily\thechapter}%
}}%
\rotatebox{90}{%
\resizebox{%
\heightof{\usebox{\feline@chapter}}+\depthof{\usebox{\feline@chapter}}}%
{!}{\scshape\so\@chapapp}}\quad%
\raisebox{\depthof{\usebox{\feline@chapter}}}{\usebox{\feline@chapter}}%
}
\newcommand\feline@chm[1][4cm]{%
\sbox\feline@chapter{\feline@chapter@marker[#1]}%
\makebox[0pt][l]{% aka \rlap
\makebox[1cm][r]{\usebox\feline@chapter}%
}}
\makechapterstyle{daleif1}{
\renewcommand\chapnamefont{\normalfont\Large\scshape\raggedleft\so}
\renewcommand\chaptitlefont{\normalfont\huge\bfseries\scshape\color{nicered}}
\renewcommand\chapternamenum{}
\renewcommand\printchaptername{}
\renewcommand\printchapternum{\null\hfill\feline@chm[2.5cm]\par}
\renewcommand\afterchapternum{\par\vskip\midchapskip}
\renewcommand\printchaptertitle[1]{\chaptitlefont\raggedleft ##1\par}
}
\makeatother
\chapterstyle{daleif1}


\begin{document}
\frontmatter
\title{MBP Merger Tree}
\author{Juhan Kim}
\affil{Center for Advanced Computation @KIAS, Seoul, Republic of Korea}
\maketitle

\chapter{Preface}
This document contains the information on the GID-sorted merger tree of the most-bound particle (MBP).
The MBP merger tree was developed for the One-to-One Correspondence (or in a 
more popular term, abundance matching) model.
The tidal disruption time scale is adopted to calculate the survival time of the satellite halos after the merger.

\mainmatter
\chapter{GID-Sorted Merger Tree}
The original Merger Tree for the MBP (most bound particle) has a caveat to nearly prevent the analysis of a huge and fat merger tree
because each merger line of MBP is nearly randomly scattered around the data file. To calculate the merger time and to find the
neighboring MBP, one needs to search all the MBP lines for the MBP's in the same halo. 
Therefore, after getting some requesting
email from Dr. S. E. Hong, I made a new MBP merger tree ordered with GID number to group together tree lines residing in the same FoF halos.

\section{Definitions}
\subsection{FoF halo}
Usually, we have used the Friend-of-Friend (FoF) scheme to identify a virialized halo from the simulation particle data. 
The FoF has been found to match very well with the analytic halo mass functions and to be well established in the cosmological context.
Therefore, from now on, whenever we mention a halo, we mean the FoF halo.


\subsection{Most Bound Particle}
The definition of the MBP (most bound particle) is the particle whose total energy is the biggest in negative value
in a halo which just forms. Therefore, a halo carries on the MBP throughout the simulation run. 
The MBP is sometimes assumed to be an ID of the halo. The MBP once created is not destroyed whenever it 
is among the member particles of some halos. 

To be more specific, for each time step, we identify the MBPs found in the previous step and check whether a halo 
contains one (or multiple) MBP(s). If a halo contains no MBP, then we calculate the total energy of the member particles,
and assign the most bound particle to the halo as an MBP. If one (multiple) MBP(s) are found, we do not have to make a new MBP.


\subsection{Attached quantities to MBP}
Strictly speaking, MBP is a particle with its own particle index, position, and peculiar velocity. But as we assume
that the MBP represents a galaxy in the halo, we can assign several halo properties to the MBP, such as
mean velocity ($v_h$), velocity dispersion ($\sigma_v$), radius ($r_{vir}$), shape ($a$, $b$, $c$)\footnote{Currently, this information is not available}, and rotation
($\lambda$ \& $\boldsymbol{L}$). In the current version, the shape, angular momentum, and spin value are
not included in the Merger Tree.


\subsection{Central \& Satellite MBP}
If two halos (with different masses) meger together and form a larger halo, then the two MBP's might be inside the merger halo.
In this case, we name the MBP coming the more massive halo the central (major) MBP and the other MBP the satellite MBP.
The satellite MBP's quantities (mass, shape, radius) 
are just set to freeze to the end of the merger tree line, except the $\sigma_v$, which
is set to the merger halo's $\sigma_v$. This means that all the MBP's in a halo are set to have the same $\sigma_v$
\footnote{In the next upgrade, we will adopt the velocity dispersion just before the first merging to a larger system
as the MBP's $\sigma_v$}.
Of course, the MBP's position and velocity are updated by tracing the particle movement inside the halo.

\subsection{GID}
GID is the global identification number stored in a $C$ format of $size\_t$. It is unique and one may search for
other MBPs' history using this GID value.

\subsection{Major GID}
The major GID is the GID number of the major MBP for each merger step. It may change with time because the
major MBP may change with time.


\section{Original Merger Tree}

The data structure of ``SavedHaloType'' is the original merger tree format, which has the following information.

\noindent\rule{\textwidth}{1pt}
$\cdots$
\begin{lstlisting}[language=c]
#ifdef XYZDBL
typedef double POSTYPE;
#else
typedef float POSTYPE;
#endif

typedef struct SavedHaloType{
    INDXTYPE mbp; /* mbp ID */
    IDTYPE nowhid,majorglobalid; 
    /* nowhid is halo id at the redshift and majorglobalid is global array 
    id to the major subhalo (or mbp ) */
    POSTYPE x,y,z,hx,hy,hz; 
    /* in unit of h^-1 Mpc in comoving coordinate (x,y,z) is position of 
    the mbp and (hx,hy,hz) is the CoM position of the hosting FoF halo */
    float mass,fofhmass; 
    /* in unit of M_sun/h 
       mass is the mass of MBP or halo (if one mbp is in the halo) just before 
       entering larger halo region.  fofhmass is the hosting halo mass. */
    float vx,vy,vz,hvx,hvy,hvz; 
    /* in unit of km/second  (vx,vy,vz) is the peculiar velocity of the mbp &
     (hvx,hvy,hvz) is the peculiar velocity of the hosting FoF halo. */
    short int statusflag, linkflag;
    /* These two members are internally used for making MBP */
} SavedHaloType;

\end{lstlisting}
$\cdots$

\noindent\rule{\textwidth}{1pt}
In a big simulation we usually turn on ``XYZDBL'' compiling option, which makes POSTYPE type a $double$ precision.
``INDXTYPE'' has been usually set to $long$ in a $C$ format.\\
IMPORTANT: note that ``majorglobalid'' in the struct SavedHaloType indicates the major MBP's GID.



\section{How to Build the Original Merger Tree}
First, we have to run a program to extract the physical properties of the FoF halos as;\\
\texttt{\% mpirun -np ** ./mpimf2.exe [step number]}\\
which produces several files including a file ``\texttt{FoFHaloQuantities.\%.5d.bin}''.\\
Then, run a code as\\
\texttt{\% mpirun -np ** ./sim2link.mpi.mod2.exe [step before] [step after]} \\
which links MBP between two adjacent steps.\\
Finally, run a code to make a chain of links as\\
\texttt{\% mpirun -np ** ./mkmerger.exe}\\
to get the final merger tree.
To sum up all the merger tree files into a single file,
one has to run \\
\texttt{\% makeone}\\
.

\section{Data Type of Sorted-MBP}

Below is the snippet of ``SMG.h'' containing the structure information of the output data file of Sorted-Merger Tree.

\noindent\rule{\textwidth}{1pt}
$\cdots$
\begin{lstlisting}[language=c]
typedef struct MajorGid{
    IDTYPE gid, newgid, mgid;
    IDTYPE majorgid[NSTEPS];
} MajorGid;

typedef struct GID{
    IDTYPE gid, newgid;
} GID;

typedef struct SigV{
    float sigv[NSTEPS];
}SigV;
typedef struct GidMergeHistory{
    IDTYPE gid;
    SavedHaloType step[NSTEPS];
    SigV sigv;
} GidMergeHistory;
\end{lstlisting}
$\cdots$

\noindent\rule{\textwidth}{1pt}

Here, note that the data type of IDTYPE is $long$ in $C$ format and NSTEPS is the number of time step of the merger tree.
The output file has a main content of array of ``structure GidMergeHistory'', which contains the GID (gid),
main MBP physical quantities in step[NSTEPS], and $\sigma_v$ in sigv (km/second).

\section{FoF Halo}
At multiple time step during the simulation run, we dumped the FoF halo data (particle position \& velocity).
And we apply the FoF algorithm to find the virialized halos. The linking length is set to be two times the mean
particle separation.
\subsection{MBP}
The MBP of a halo is found by calculating the total energy $E$ as
\begin{equation}
E = {1\over2} m \left({d \boldsymbol{r} \over dt}\right)^2 + \Phi(\boldsymbol{r})
\end{equation}
where $\boldsymbol{v}$ is the sum of peculiar velocity and Hubble flow.
The potential $\Phi$ can be obtained from
\begin{equation}
\Phi(r) \equiv m\sum_i {Gm_i\over \sqrt{ (\boldsymbol{r}_i - \boldsymbol{r})^2 + \epsilon^2}},
\end{equation}
where $\epsilon$ is the force resolution of the simulation and $i$ runs for every member particle of the halo.
The particle with the lowest value of $E$ (or most negative total energy) are named the MBP of the halo.

\subsection{CoM of a Halo}
The Center of Mass (CoM) of a halo is calculated as the mean position of member particles. Also the group velocity
is the average of the peculiar velocity of them.

\subsection{Shape of a Halo}
The shape of a halo is measured by finding the eigen values and vectors of the shape tensor as
\begin{equation}
\boldsymbol{T}_{ij} \equiv \sum_i (\boldsymbol{r}_i \boldsymbol{r}_j)
\end{equation}
where $\boldsymbol{r}_i$ is the relative distance of a particle from the CoM of the halo.
The ratios of eigen values are the ratios of the lengths of the axes.



\chapter{Merger Tree of MBP}
\section{Basic Concept}

\subsection{Merger Events}
The basic concept in building the MBP merger tree is to list the merger history of one MBP in a single line
and the merger events can possibly be marked by the change of major GID (if it is a satellite MBP).
When a halo just forms, the halo has one MBP and I allocate a unique GID number to it. 
Also the major GID of the MBP is same as its GID in this case.
But when there happens a merger event, the major GID changes to the most massive merging counterpart's GID.
\subsection{Mass of MBP}
The mass of a MBP is the mass of its hosting halo when there is only one MBP or when the MBP is the major MBP
in the halo. If a smaller halo with one/multiple MBP('s) mergers to form a larger halo, the major MBP of the samller halo
has the halo mass measured just before the merger event. Therefore, all the satellite halos (it doesnot matter whether they
are satellites in satellites) contains their mass when they are merged to larger halos previously.
\subsection{fofhmass}
``fofhmass'' contains the halo mass of the major MBP or hosting halo mass
identified by the FoF with $l_{ linking} = 0.2\times$ the mean particle separation. Usually, I set the minimum
halo mass $m_{h}^c = 30\times m_p$, where $m_p$ is the particle mass.

\subsection{mbp, nowhid, \& majorglobalid}
``mbp'' is the particle index of the MBP.
``nowhid'' is the halo id in the FoF halo catalog.
``majorglobalid'' is the major GID.

\section{How To Sort}
We sort the merger line of mbp according to their major GID at $z=0$,
which means all the tree lines below a major MBP at $z=0$ are grouped together in the file.
This may help to calculate the merger effects. 

\section{Valuable Notes}
\subsection{Disappearance of MBP}
In some cases, the MBP of a halo may be kicked off from its host halo for some reason and goes somewhere not in other halos
(maybe expelled to the intergalactic medium?). 
In this case the MBP line stops there
(of course the MBP's variables are set to null in the following steps). Then, a new MBP is selected in the halo and a new MBP line 
starts thereafter.
\subsection{Flucutating of Small Halos}
When a halo just forms, its mass is just above some critical (input) mass. But in the next time step the halo may lose the mass
by some physical cause or numerical fluctuations. In this case, it may generate multiple MBP's for a sinlge physical halo.
Of course, only one MBP may evolve to the subsequent time steps and the other MBP's does not contain any information
after the subsequent steps.

\subsection{NULL VALUES}
If there is a null/0/negative value in ``mbp'', it may indicate that the MBP is not found at the epoch.

\subsection{Mass Problem}
Sometimes, the total mass sum of halos before the merger and after merger may be different.
Some mass loss can happen and analysis should reflect this mass loss during the merger.
\subsection{}



\chapter{Codes}
\section{Example Code To Read Data}
The example code to read the data file can be found
in \\
$\//scratch:
\//scratch\//kjhan\//Nbody\//WMAP5\//2LPT2048p1024s\//Om0.26w-1\\ 
\//FoFPSB\//
L0.2\//Mass\//Merging\//MKMerger$,
\noindent where you can see a source code with name ``checkSGMT.c''.


\noindent\rule{\textwidth}{1pt}
\begin{lstlisting}[language=c]

/* icc -o checkSGMT checkSGMT.c -DINDEX -DXYZDBL -g */
#include<stdio.h>
#include<stdlib.h>
#include<math.h>
#include<stddef.h>
#include "merger.h"
#include "mkmerger.h"
#include "SMG.h"
#define NSTEPS 131 /* number of merger steps */

int *nsteps;
float *reds;
float boxsize, hubble, npower, omep, omeplam, bias,astep,amax;
float anow,omepb;
int nx, nspace;
int msteps,*nsteps;
int main(int argc, char *argv[]){
    IDTYPE i,j,k;
    GidMergeHistory *bp;
    IDTYPE np;

    FILE *fp = fopen(argv[1],"r");
    fread(&boxsize, sizeof(float), 1,fp);/* simulation box size in h^-1 Mpc */
    fread(&hubble, sizeof(float), 1,fp); /* hubble parameter, for example, h=0.72 */
    fread(&npower, sizeof(float), 1,fp); /* power spectral index */
    fread(&omep, sizeof(float), 1,fp); /* omega_matter at present */
    fread(&omepb, sizeof(float), 1,fp); /* Omega_baryon0 at present */
    fread(&omeplam, sizeof(float), 1,fp); /* Omega_lambda at present */
    fread(&bias, sizeof(float), 1,fp); /* bias factor or inverse of sigma_8 */
    fread(&astep, sizeof(float), 1,fp); /* simulation astep */
    fread(&msteps, sizeof(float), 1,fp); /* number of steps in the merging tree */
    nsteps = (int*)malloc(sizeof(int)*msteps);
    reds = (float*)malloc(sizeof(float)*msteps);
    fread(nsteps, sizeof(int), msteps,fp); /* step number save to merging tree */
    fread(reds, sizeof(float), msteps,fp); /* redshift info. save to merging tree */


    bp = (GidMergeHistory*)malloc(sizeof(GidMergeHistory)*10000);
    while((np = fread(bp,sizeof(GidMergeHistory),10000,fp))>0){
        printf("%ld\n",bp[0].gid);
    }
    free(bp);
}


\end{lstlisting}
\noindent\rule{\textwidth}{1pt}
The data file has a header containing basic information of the simulation such as the simulation box size,
Hubble parameter, $\Omega_m$, $\Omega_\Lambda$, $b_8$ and so on. Also it has the number of Merger Tree steps ($msteps$)
with a corresponding redshift. Below the header, it attaches the main data of ``GidMergeHistory''.
The input data file name is ``SortByMajorGidMergingTree.00130.dat'' which is in \\
$\//scratch:
\//scratch\//kjhan\//Nbody\//WMAP5\//2LPT2048p1024s\//Om0.26w-1 
\//FoFPSB\//
L0.2$.

\noindent One may compile this code as\\
\texttt{\% icc -o checkSGMT checkSGMT.c -DINDEX -DXYZDBL -g}\\
and you can run the code in debugging mode as

\noindent\rule{\textwidth}{1pt}
\begin{lstlisting}
[kjhan@annapurna L0.2]% gdb checkSGMT
GNU gdb (GDB) CentOS (7.0.1-45.el5.centos)
Copyright (C) 2009 Free Software Foundation, Inc.
License GPLv3+: GNU GPL version 3 or later <http://gnu.org/licenses/gpl.html>
This is free software: you are free to change and redistribute it.
There is NO WARRANTY, to the extent permitted by law.  Type "show copying"
and "show warranty" for details.
This GDB was configured as "x86_64-redhat-linux-gnu".
For bug reporting instructions, please see:
<http://www.gnu.org/software/gdb/bugs/>...
Reading symbols from ./checkSGMT...(no debugging symbols found)...done.
(gdb) l
16      int nx, nspace;
17      int msteps,*nsteps;
18
19
20      int main(int argc, char *argv[]){
21              IDTYPE i,j,k;
22              GidMergeHistory *bp;
23              IDTYPE np;
24
25              FILE *fp = fopen(argv[1],"r");
(gdb) l
26              fread(&boxsize, sizeof(float), 1,fp);/* simulation box size in h^-1 Mpc */
27              fread(&hubble, sizeof(float), 1,fp); /* hubble parameter, for example, h=0.72 */
28              fread(&npower, sizeof(float), 1,fp); /* power spectral index */
29              fread(&omep, sizeof(float), 1,fp); /* omega_matter at present */
30              fread(&omepb, sizeof(float), 1,fp); /* Omega_baryon0 at present */
31              fread(&omeplam, sizeof(float), 1,fp); /* Omega_lambda at present */
32              fread(&bias, sizeof(float), 1,fp); /* bias factor or inverse of sigma_8 */
33              fread(&astep, sizeof(float), 1,fp); /* simulation astep */
34              fread(&msteps, sizeof(float), 1,fp); /* number of steps in the merging tree */
35              nsteps = (int*)malloc(sizeof(int)*msteps);
(gdb) l
36              reds = (float*)malloc(sizeof(float)*msteps);
37              fread(nsteps, sizeof(int), msteps,fp); /* step number save to merging tree */
38              fread(reds, sizeof(float), msteps,fp); /* redshift info. save to merging tree */
39
40
41              bp = (GidMergeHistory*)malloc(sizeof(GidMergeHistory)*10000);
42              while((np = fread(bp,sizeof(GidMergeHistory),10000,fp))>0){
43                      printf("%ld\n",bp[0].gid);
44              }
45              free(bp);
(gdb) break 43
Breakpoint 1 at 0x4008b8: file checkSGMT.c, line 43.
(gdb) run SortByMajorGidMergingTree.00130.dat
Starting program: checkSGMT SortByMajorGidMergingTree.00130.dat
warning: no loadable sections found in added symbol-file system-supplied DSO at 0x2aaaaaaab000

Breakpoint 1, main (argc=2, argv=0x7fffffffd5b8) at checkSGMT.c:43
43                      printf("%ld\n",bp[0].gid);
(gdb) print bp[0].gid
$1 = 0
(gdb)  print bp[0].sigv
$2 = {sigv = {200.545746, 229.547607, 240.176117, 261.648102, 320.950165, 412.169983, ............ }}
(gdb) 
\end{lstlisting}

\noindent\rule{\textwidth}{1pt}
This is a typical example run with data on 2LPT2048p1024s.


\section{Source Codes Building Sorted Merger Tree}
I build two steps to make sorted GID merger tree.
First assign a new GID to the original merger tree data and, then, sort them according to the new GID.

\begin{itemize}
\item How to assign a new GID:

Source codes:\\
\noindent\rule{\textwidth}{1pt}
\begin{lstlisting}
Make.SMB  /* make file */
SortMajorGid.c /* main source code */
pqsort.SMG.c   /* needed for parallel sorting */
SMG.h
pqsort.SMG.h
\end{lstlisting}
\noindent\rule{\textwidth}{1pt}
To run the executable, type as\\
\texttt{\% mpirun -np ** ./SortMajorGid.exe [redshift]}\\
where redshift is the redshift to sort GID.
One has to note that the internal structure can change depending on the final redshift (it could be $z=0$ or some
other redshift). Please check ``iredshift'', ``NSTEPS'', and ``MSTEPS''.

\item How to sort and dump the data:
Source codes:\\
\noindent\rule{\textwidth}{1pt} 
\begin{lstlisting}
Make.SSMB  /* make file */
Sort_SortMajorGid.c /* main source code */
pqsort.SMG.c   /* needed for parallel sorting */
SMG.h
pqsort.SMG.h
\end{lstlisting}
\noindent\rule{\textwidth}{1pt}
To run the executable, type as\\
\texttt{\% mpirun -np ** ./Sort\_SortMajorGid.exe [input file header] [no. of input files]}\\


\end{itemize}

Any comment or suggestion will be welcome.



\end{document}


